% Preguntas conceptuales. Escribe cada respuesta dentro de \respuesta{...}.

\subsection*{C1. Errores y teorema de Lax--Richtmyer\hfill[2 pts]}
\addcontentsline{toc}{subsection}{C1. Errores y Lax--Richtmyer}
Distinga error de truncamiento, error local y error global. Enuncie la relación entre
consistencia, estabilidad y convergencia (Lax--Richtmyer).
\respuesta{}

\subsection*{C2. Del error local al error global\hfill[2 pts]}
\addcontentsline{toc}{subsection}{C2. Del error local al global}
Explique por qué un método de un paso con error local $\Order{h^{p+1}}$ tiene error global
$\Order{h^p}$ y qué papel juega la estabilidad.
\respuesta{}

\subsection*{C3. Radio espectral y razón de dominancia\hfill[2 pts]}
\addcontentsline{toc}{subsection}{C3. Radio espectral y razón de dominancia}
Defina el radio espectral y la razón de dominancia. ¿Cómo determinan la rapidez de convergencia
del método de las potencias y de una iteración $\vect{x}^{(n+1)} = \mat{M}\vect{x}^{(n)} + \vect{c}$?
¿Qué consecuencia tiene un $\rho_D$ cercano a 1 en un núcleo grande y débilmente acoplado?
\respuesta{}

\subsection*{C4. Condiciones de frontera en difusión neutrónica\hfill[2 pts]}
\addcontentsline{toc}{subsection}{C4. Condiciones de frontera}
Interprete físicamente las condiciones de Dirichlet, Neumann y Robin en difusión neutrónica
(vacío extrapolado, frontera reflejante, albedo) y describa su efecto en la matriz del sistema.
\respuesta{}

\subsection*{C5. Esquemas para difusión transitoria\hfill[2 pts]}
\addcontentsline{toc}{subsection}{C5. Esquemas para difusión transitoria}
Compare Euler explícito, Euler implícito y CN para difusión transitoria en cuanto a estabilidad,
orden y oscilaciones espurias. ¿Por qué se factoriza una sola vez la matriz si $\Delta t$ es constante?
\respuesta{}
